% Optional insertion before the existing kernel subsection in manuscript/draft.
% Analysis 018, hand-edited descriptive results fragment, 2026-09-08.
% Evidence: figure_data.json and observations/O001-O007 in this analysis.
% Generators: evidence.py, plots.py, 01_build.py. No manuscript file was edited.
% Requirements: graphicx, booktabs, amsmath (already present in draft/main.tex).
% Override this path when compiling from a different working directory.
\providecommand{\resultsassets}{../../analyses/018-2026-09-08-results-materials}

\subsection{Quality--sparsity trade-offs at 14M}
\label{sec:results-overview}

% Evidence: observations/O001-overview.md, O006-posthoc-completeness.md.
The detailed study contains 30 conditions from the main intervention ladder
and five historical A4 conditions with OL1 applied only at $h$, denoted
A4-OL1[$h$]. All share initialization and data-order seed 1234, 712 optimizer
updates and 1,493,172,224 input tokens. Training uses FP16 dynamic loss scaling,
FP32 parameters and AdamW with global batch 1,024; the peak learning rate is
$10^{-3}$. Each endpoint pairs loss and model-wide sparsity in the same complete
eager validation pass over 338 blocks, excluding the 1,444-token tail.

Figure~\ref{fig:results-overview} spans distinct trade-off regimes. Local L1
pressure at the ReLU FFN hidden activation yields the lowest observed loss:
at $\lambda=1$, loss is 5.1023 with 3.9493\% model-wide sparsity, compared with
5.2696 and 2.7141\% for unpressured A1-H. Broader placement reaches a different
region: A7-OL1 at $\kappa=0.5$ reaches 27.4827\% at loss 5.8294.
The historical A4-OL1[$h$] condition at $\kappa=0.05$ reaches 9.0356\% at
loss 5.1956. It is distinct from corrected four-site A4-OL1, which contributes
no globally nondominated point to the main 30-condition trained ladder.

\begin{figure}[t]
\centering
\includegraphics[width=\linewidth]{\resultsassets/figures/01-14m-overview.pdf}
\caption{Complete 14M trained overview with the 20 A0/A1-H uniform-clipping
controls. Colors identify recipes; gray dotted curves are evaluation-only
clipping. Rings mark nondominated trained points, with solid and dotted
envelopes for the main ladder and the pool including historical A4-OL1[$h$].
The right panel retains the full loss range, including eight clipping points
above the left panel. Connecting lines guide the eye through evaluated points;
they do not establish attainable intermediate operating points.}
\label{fig:results-overview}
\end{figure}

The complete retained post-hoc comparison additionally contains 150 evaluations
across 15 source checkpoints. Clipping the $\lambda=1$ L1 checkpoint to target
$p=0.1$ gives 4.7982\% sparsity at loss 5.1070. These are uniform TEAL-style
targets, not greedy allocation, and the retained clipping pool does not cover
A7 or corrected A4-OL1. Thus the available comparison does not establish the
best possible training-plus-clipping composition.

\subsection{Pressure's added value depends on threshold and placement}
\label{sec:results-effects}

% Evidence: observations/O002-blocked-effects.md; tables/pressure-and-gate-effects.tex.
Figure~\ref{fig:results-effects} shows child-minus-parent effects with matched
thresholds. Adding OL1 to A4 improves both outcomes at $\kappa=0$ and $0.01$.
At $\kappa=0.5$, it adds 2.4979 sparsity points at a 0.3783 loss cost.
For A7, OL1 instead slightly worsens both outcomes at zero threshold, adds
1.3717 points for a 0.000819 loss increase at $\kappa=0.1$, and adds 12.0959
points for a 0.1265 loss increase at $\kappa=0.5$. Pressure therefore has no
uniform benefit over the tested thresholds.

\begin{figure}[t]
\centering
\includegraphics[width=\linewidth]{\resultsassets/figures/02-blocked-intervention-effects.pdf}
\caption{Blocked matched contrasts at 14M. Each x position aligns the change
in loss (top) with the change in model-wide sparsity (bottom). Negative loss
and positive sparsity differences are favorable. Every block names its own
parent and child; $\lambda$ and $\kappa$ are treatment levels, not statistical
replicates. The sequence organizes separately trained conditions and is not a
cumulative training curriculum.}
\label{fig:results-effects}
\end{figure}

At fixed topology these contrasts isolate addition of the specified pressure
objective. Across A4 and A7, however, both pressure sites and their equal-tensor
normalization change. The response difference is therefore recipe dependent;
without A7 gates with four-site pressure it cannot identify a fixed-objective
gate-by-pressure interaction. Similarly, expanding historical pressure from
$h$ to four A4 sites raises loss by 0.2425--0.3196 at every matched threshold,
but does not identify which added site causes that cost.
At zero threshold the additional unpressured A7 gates are identities; the
small measured A4/A7 residual is a numerical/implementation control discrepancy,
not a benefit of an identity gate. These one-seed comparisons do not support
statistical significance or equivalence claims.

\subsection{Selected recipes transfer, with a changed low-threshold response}
\label{sec:results-transfer}

% Evidence: observations/O003-scale-transfer.md and O004-operation-accounting.md.
At $\kappa=0.5$, A7-OL1 improves on A4-OL1 at every tested size: model-wide
sparsity increases by 14.7692, 5.0057 and 9.0241 points at 14M, 70M and 410M,
while loss decreases by 0.2086, 0.1736 and 0.0703, respectively.
At zero threshold, A4-OL1 is better on both outcomes at 14M and 70M, whereas
A7-OL1 is better at 410M. This is partial recipe transfer rather than a uniform
scaling pattern. The larger cohorts lack unpressured A4/A7, so they cannot
replicate the isolated pressure contrasts.

\begin{figure}[t]
\centering
\includegraphics[width=\linewidth]{\resultsassets/figures/03-scale-transfer-and-ceilings.pdf}
\caption{Selected recipes across scales. Loss is relative to the same-size A0.
The top row uses raw model-wide sparsity, with topology ceilings marked by
colored dotted lines. The bottom row divides by each recipe's own ceiling;
the dotted line marks a ratio of one. All five doses are connected in order,
with $0$ and $0.5$ labeled. A0's normalized value is undefined. The ratio
includes natural zeros outside selected reach and is not generally bounded
by 100\%; different recipes also have different reachable workloads.}
\label{fig:results-transfer}
\end{figure}

The A7 reach ceiling rises from 29.952\% to 49.424\% and 87.245\%, while
its high-threshold endpoints realize 91.75\%, 82.15\% and 92.40\% of those
ceilings. This normalization exposes the contribution of architectural reach
to raw sparsity growth, but is not a size-invariant performance score.
Indeed, A4-OL1 realizes a larger fraction of its narrower ceiling at every
matched dose and size, including where A7-OL1 is better on both raw outcomes.
Operation-level accounting separates each operation's learned zero-product
rate from its share of the full model workload.

All sizes receive the same token budget, giving 106.1, 21.2 and 3.7
tokens per parameter. The 410M peak learning rate is $3\times10^{-4}$,
compared with $10^{-3}$ at the two smaller sizes. One seed per size and
unequal exposure per parameter limit extrapolation; the observed reversal
does not identify undertraining or a particular optimizer mechanism.

\subsection{Small attention activations and exact-zero attention operands}
\label{sec:results-activations}

% Evidence: observations/O005-activation-case-study.md; tables/activation-statistics.md.
At 14M and $\kappa=0.5$, A4-OL1 and A7-OL1 both have more than 99.8\%
exact zeros at $h$, yet differ strongly in attention. Under A4-OL1,
47.64\% of queries and 56.12\% of keys have magnitude at most 0.01, while
their exact-zero fractions remain below 0.001\%. Under A7-OL1, queries and
keys are 93.54\% and 94.54\% exactly zero, and values are 98.71\% zero.
Query RMS is smaller under A4-OL1 (0.146 versus 0.483), showing why magnitude
reduction alone does not identify the model-wide zero-product opportunity.

\begin{figure}[t]
\centering
\includegraphics[width=\linewidth]{\resultsassets/figures/05-activation-mass-grid.pdf}
\caption{Retained activation magnitude bands and RMS at 14M. Rows compare
$\kappa=0$ and $0.5$ for A4-OL1/A7-OL1 against the same repeated A0 baseline.
Columns use six common diagnostic sites; attention output is after $W_o$,
not the pre-projection gated context $z$. Four disjoint empirical magnitude
bands are count-pooled over all layers and validation blocks. Lines connect
categorical bands; they are not density estimates. The y scale is logarithmic
above 0.01\% and linear near zero.}
\label{fig:results-activations}
\end{figure}

These retained diagnostics contain exact/near-zero counts and moments rather
than full histograms. The case study identifies how the recipes' observed
zeros differ by site, without establishing a causal spillover mechanism.
The following kernel case study tests whether an implementation can exploit
the corresponding opportunity on a declared hardware workload.
